\chapter{Krull Topology and the Inverse Limit}\label{ch:krull-topology}

Let \(L/K\) be a Galois extension, and put
\(G=\operatorname{Gal}(L/K)\). Products in \(G\) are compositions,
with the right factor applied first. For an intermediate field \(E\),
write \(G_E=\operatorname{Gal}(L/E)\), viewed as a subgroup of \(G\).
For a finite subset \(S\subseteq L\), put
\[
 G(S)=\{\sigma\in G:\sigma(x)=x\text{ for all }x\in S\}.
\]
In particular, \(G(S)=G_{K(S)}\). Fixing \(S\) fixes every expression
obtained from its elements and \(K\) by the field operations, and hence
fixes the field they generate.

\section{The neighborhood basis and restriction}

\begin{proposition}[Finite pointwise stabilizers]\label{res:T01}
There is a unique topology with basis consisting of the left cosets
\(\sigma G(S)\), where \(S\) ranges over the finite subsets of \(L\).
In this topology \(G\) is a topological group, and the groups \(G(S)\)
are an open neighborhood basis at the identity. This is the
\emph{Krull topology}.
\end{proposition}

\begin{proof}
The subsets \(G(S)\) are subgroups, with
\[
 G(S)\cap G(S')=G(S\cup S'),\qquad
 \sigma G(S)\sigma^{-1}=G(\sigma S).
\]
For the second identity, \(\sigma h\sigma^{-1}\) fixes each \(\sigma x\)
if and only if \(h\) fixes each \(x\). Every point belongs to a proposed
basic open set. If \(\gamma\) belongs to two basic cosets, each coset can
be written using \(\gamma\) as representative. Their intersection then
equals \(\gamma G(S\cup S')\). This proves the basis property.
Equivalently, \(\sigma G(S)\) consists of all automorphisms which agree
with \(\sigma\) on \(S\).

For continuity of multiplication, fix a basic neighborhood
\(\sigma\tau G(S)\) at a product. If \(\tau'\) agrees with \(\tau\) on
\(S\) and \(\sigma'\) agrees with \(\sigma\) on \(\tau S\), then
\(\sigma'\tau'(x)=\sigma\tau(x)\) for \(x\in S\). Thus the product of
the neighborhoods \(\sigma G(\tau S)\) and \(\tau G(S)\) is contained
in the given neighborhood. For inversion, if \(\sigma'\) agrees
with \(\sigma\) on \(\sigma^{-1}S\), then
\(\sigma'(\sigma^{-1}x)=x\) for \(x\in S\). Applying the inverse of
\(\sigma'\) proves that \((\sigma')^{-1}\) agrees with \(\sigma^{-1}\)
on \(S\). This supplies the required basic neighborhood at \(\sigma\).
In a topological group translations are homeomorphisms, so the identity
basis determines the entire topology. See
\cite[Section 7, following 7.2, and Proposition 7.6]{milne}.
\end{proof}

Let \(\mathcal E\) be the set of finite Galois intermediate fields of
\(L/K\). For \(E\in\mathcal E\), put
\[
 Q_E=\operatorname{Gal}(E/K),\qquad
 r_E\colon G\longrightarrow Q_E,\quad \sigma\longmapsto \sigma|_E,
 \qquad N_E=\ker r_E=G_E.
\]
We regard every \(Q_E\) as a discrete group. Theorem~\ref{res:F04}
gives the two indexing properties used here: \(\mathcal E\) is
directed by inclusion, and its union is \(L\).

\begin{proposition}[A basis of finite quotient kernels]\label{res:T02}
The subgroups \(N_E\), for \(E\in\mathcal E\), form an open neighborhood
basis at the identity. Each is clopen and normal of finite index, and
\(\bigcap_{E\in\mathcal E}N_E=\{1\}\). If \(L/K\) is finite, the Krull
topology is discrete.
\end{proposition}

\begin{proof}
For each \(E\in\mathcal E\), choose a \(K\)-basis of \(E\) and let \(S_E\)
be its finite set of elements. Then \(N_E=G(S_E)\), since a
\(K\)-automorphism fixing the basis pointwise fixes \(E\). Conversely,
Proposition~\ref{res:F03} places any finite \(S\subseteq L\) in some
\(E\in\mathcal E\). For that \(E\) one has \(N_E\subseteq G(S)\), giving
the claimed basis. Proposition~\ref{res:F07} identifies \(N_E\) as the
kernel of the surjection \(r_E\), so it is normal and its index is
\(|Q_E|\). Every open subgroup has open cosets. The complement of
\(N_E\) is the union of its other cosets, making \(N_E\) closed as well.
An automorphism in all \(N_E\) fixes their fields' union \(L\) pointwise.
If \(L/K\) is finite we may take \(E=L\), for which \(N_E=\{1\}\).
Translations then give open singletons throughout \(G\).
Compare \cite[Propositions 7.6--7.8]{milne} and
\cite[Section 2.11]{rz}.
\end{proof}

\begin{proposition}[Initial topologies and continuous action]\label{res:T03}
The Krull topology is the initial topology of the evaluations
\(\sigma\mapsto\sigma(x)\) to discrete \(L\), for all \(x\in L\). It is
also the initial topology of the restrictions \(r_E\) to discrete
\(Q_E\), for all \(E\in\mathcal E\). The action \(G\times L\to L\) is
jointly continuous when \(L\) is discrete. If \(M/K\) is a Galois
intermediate extension, then restriction
\(G\to\operatorname{Gal}(M/K)\) is continuous for the Krull topologies.
\end{proposition}

\begin{proof}
The nonempty inverse image of a singleton under the evaluation at \(x\)
is a coset of \(G(\{x\})\), and any \(\sigma G(S)\) is the intersection
of the corresponding fibers for \(x\in S\). These identities give the
first topology. Each nonempty fiber of \(r_E\) is a coset of \(N_E\).
Conversely all these cosets are fibers. By Proposition~\ref{res:T02}
they form a basis, giving the second initial topology.

At \((\sigma,x)\), the action is constant on
\(\sigma G(\{x\})\times\{x\}\). This is an open neighborhood, which
proves joint continuity. It also yields the corresponding universal
property for an action map: if a map \(u\colon X\to G\) from a
topological space has continuous induced map
\((z,x)\mapsto u(z)(x)\) on \(X\times L\), then each
\(z\mapsto u(z)(x)\) is continuous by fixing \(x\), so \(u\) is
continuous by the first initial description.

For the restriction to \(M\), Proposition~\ref{res:F07} gives the
homomorphism, since \(M/K\) is Galois. For \(x\in M\), its composite
with evaluation at \(x\) is the evaluation in \(L\) with values in the
discrete subset \(M\). This composite is continuous as a map into
discrete \(M\), because the inverse images of its singletons are
the same as the corresponding inverse images for evaluation in \(L\).
Applying the first initial description to \(M/K\) gives continuity
of restriction, including when \(M/K\) is infinite.
See \cite[Propositions 7.6--7.7]{milne} and
\cite[Tags 0BMJ and 0BMK]{stacks}.
\end{proof}

\section{An inverse system of finite groups}

\begin{proposition}[The compatible-tuple limit]\label{res:T04}
For \(E\subseteq E'\) in \(\mathcal E\), let
\(r_{E',E}\colon Q_{E'}\to Q_E\) be restriction. These onto maps
form an inverse system. With coordinatewise group operations and
the subspace topology from \(\prod_{E\in\mathcal E}Q_E\), the group
\[
 P=\left\{(\sigma_E)_{E\in\mathcal E} :
       r_{E',E}(\sigma_{E'})=\sigma_E
       \text{ whenever }E\subseteq E'\right\}
\]
is its inverse limit in topological groups.
\end{proposition}

\begin{proof}
The restrictions make sense because \(E/K\) is Galois, and
Proposition~\ref{res:F07}, applied to \(E'/K\), gives surjectivity.
Restrictions are the identity when the two fields agree, and
\(r_{E'',E}=r_{E',E}r_{E'',E'}\) for \(E\subseteq E'\subseteq E''\).
The maps are continuous between discrete groups. Every compatibility
equation is preserved by multiplication and inversion because its
transition map is a homomorphism. Thus \(P\) is a subgroup of the
product, where multiplication and inversion are continuous by the
product topology.

Let \(X\) be a topological group with continuous homomorphisms
\(f_E\colon X\to Q_E\) satisfying \(r_{E',E}f_{E'}=f_E\).
Their componentwise pairing takes values in \(P\). It is a continuous
homomorphism to \(P\) by the product and subspace topologies, with all
the prescribed components. An element of \(P\) is its tuple of
components, so these conditions determine the factorization uniquely.
This is the stated limit property.
See \cite[Examples 7.22--7.23 and 7.26]{milne}.
\end{proof}

\begin{theorem}[Gluing compatible automorphisms]\label{res:T05}
The homomorphism
\[
 \rho\colon G\longrightarrow P,\qquad
 \sigma\longmapsto (\sigma|_E)_{E\in\mathcal E}
\]
is an isomorphism of groups. Its composite with projection to \(Q_E\)
is \(r_E\), and each such projection is surjective.
\end{theorem}

\begin{proof}
The finite restrictions of an automorphism are compatible by their
composition law. Two automorphisms with the same restrictions agree
on the union \(L\), proving injectivity.

Given \((\sigma_E)\in P\) and \(x\in L\), choose \(E\in\mathcal E\)
containing \(x\) and set \(\sigma(x)=\sigma_E(x)\). If \(E'\) also
contains \(x\), choose a common upper field \(E''\) by directedness.
The two proposed values equal \(\sigma_{E''}(x)\) by compatibility.
To test addition or multiplication, choose one field containing
both inputs. It also contains their sum and product, and the
corresponding \(\sigma_E\) respects those operations. The glued map
fixes \(K\) and preserves \(1\) for the same reason.

The tuple \((\sigma_E^{-1})\) is compatible because restrictions are
homomorphisms. Its glued map is a two-sided inverse to \(\sigma\):
on any \(E\), both maps preserve \(E\) and their restrictions are
inverse automorphisms. Thus \(\sigma\in G\) and \(\rho(\sigma)\) is
the prescribed tuple. The identification of projections follows
from the definition of \(\rho\), and their surjectivity follows
from Proposition~\ref{res:F07}.
See \cite[Example 7.26]{milne},
\cite[Theorem 2.11.1]{rz}, and \cite[Tag 0BU2]{stacks}.
\end{proof}

\begin{theorem}[Topological inverse limit and profiniteness]\label{res:T06}
The isomorphism \(\rho\) is a homeomorphism. Consequently \(G\) is
compact, Hausdorff, and totally disconnected, and is profinite.
\end{theorem}

\begin{proof}
Continuity of every \(r_E\) makes \(\rho\) continuous into the product
and hence into \(P\). Its inverse is continuous because every
\(r_E\) composed with that inverse is the coordinate projection
on \(P\), and the \(r_E\) generate the Krull topology by
Proposition~\ref{res:T03}.

Each \(Q_E\) is finite discrete, hence compact and Hausdorff.
Their product is compact by the Tychonoff product theorem proved in
\cite[Chapter~5, \texttt{d10:05:tychonoff}]{d10}. It is Hausdorff
because distinct tuples differ in a coordinate.
For each \(E\subseteq E'\) the compatibility equation defining \(P\)
is the inverse image of the diagonal of \(Q_E\times Q_E\) under
\[
  (\sigma_D)_D \longmapsto
      \bigl(r_{E',E}(\sigma_{E'}),\sigma_E\bigr).
\]
This map is continuous and the diagonal is closed. Taking the
intersection of these closed sets shows that \(P\) is a closed
subspace of the compact Hausdorff product. If two points of \(P\)
differ, a fiber of one coordinate separates them by a clopen
subset of \(P\). No connected subset can contain both points.
Thus \(P\) and \(G\) are totally disconnected. The presentation
by finite discrete groups gives profiniteness directly.
See \cite[Propositions 7.8 and 7.25 and Example 7.26]{milne}
and \cite[Theorem 2.11.1]{rz}.
\end{proof}

\section{Closure, cofinality, and a change of base}

\begin{proposition}[Closure on finite levels]\label{res:T07}
For any subset \(A\subseteq G\),
\[
 \overline A=\bigcap_{E\in\mathcal E}r_E^{-1}(r_E(A)).
\]
If \(\mathcal E_0\subseteq\mathcal E\) is cofinal for inclusion,
then it suffices to use \(E\in\mathcal E_0\) in this formula.
For the same cofinal subset the coordinate restriction
\(P\to\varprojlim_{D\in\mathcal E_0}Q_D\) is an isomorphism of
topological groups.
\end{proposition}

\begin{proof}
For each \(g\in G\), its basic neighborhood \(gN_E\) is
\(r_E^{-1}(\{r_E(g)\})\). It meets \(A\) if and only if
\(r_E(g)\in r_E(A)\). Membership in \(\overline A\) is equivalent
to this meeting condition for every basic neighborhood, proving the
formula. If \(E\subseteq D\), then \(N_D\subseteq N_E\), so the
cofinal subgroups still furnish a neighborhood basis. The same
proof gives the restricted closure formula.

To see the limit comparison, start with a compatible tuple over
\(\mathcal E_0\). For any \(E\in\mathcal E\), choose
\(D\in\mathcal E_0\) with \(E\subseteq D\) and define its \(E\)-entry
to be \(r_{D,E}(\sigma_D)\). If two choices \(D,D'\) are made, first
take a common upper bound in \(\mathcal E\) and then an upper
bound for it in \(\mathcal E_0\). The compatibility equations
through this upper field make the two entries equal. For
\(E\subseteq E'\), choose \(D\) containing \(E'\), so the transition
law also makes the reconstructed entries compatible. This proves
that reconstruction is inverse to coordinate restriction.
It is continuous since each reconstructed entry is the composite
of a projection on \(\mathcal E_0\) and a transition. Both maps
respect the group operations coordinatewise. Compare
\cite[Proposition 7.12(b) and Example 7.26]{milne}.
\end{proof}

\begin{proposition}[The Krull topology over an intermediate field]
\label{res:T08}
For any intermediate field \(M\) of \(L/K\), the topology on \(G_M\)
induced from \(G\) equals its intrinsic Krull topology for \(L/M\).
\end{proposition}

\begin{proof}
The extension \(L/M\) is Galois by Proposition~\ref{res:F02}. Its
Krull topology has identity neighborhoods given by the pointwise
stabilizers of finite \(S\subseteq L\) in \(G_M\), by
Proposition~\ref{res:T01} over the new base. These are precisely
\(G_M\cap G(S)\). The intersections also form the induced
neighborhood basis at the identity, and their translates by
points of \(G_M\) therefore give both topologies. This comparison
uses the same elements of the top field \(L\) on the two sides and
applies to every intermediate field \(M\).
See \cite[Theorem 2.11.3]{rz}.
\end{proof}
